\documentclass[11pt]{article}
\usepackage[margin=1in]{geometry}
\usepackage{amsmath,amssymb,amsthm,mathtools}
\usepackage{graphicx}
\usepackage{microtype}
\usepackage[hidelinks]{hyperref}

\newtheorem{theorem}{Theorem}
\newtheorem{lemma}[theorem]{Lemma}
\newtheorem{corollary}[theorem]{Corollary}
\theoremstyle{remark}
\newtheorem*{remark}{Remark}

\newcommand{\R}{\mathbb R}
\newcommand{\B}{\mathcal B}
\newcommand{\Hh}{\mathcal H}
\newcommand{\Wdot}{\dot W}
\newcommand{\norm}[1]{\lVert #1\rVert}
\newcommand{\cof}{\operatorname{cof}}
\newcommand{\supp}{\operatorname{supp}}
\newcommand{\BMO}{\operatorname{BMO}}

\title{Nowhere-Surjective Linearization of the Global Jacobian\\
and the Failure of Locally Lipschitz Sections}
\author{Partial-result packet for arXiv:1611.02223}
\date{11 August 2026}

\begin{document}
\maketitle

\begin{center}
\fbox{\parbox{0.91\textwidth}{
\textbf{Status: candidate substantial partial result, likely valid; human
review requested.}  For every $n\ge2$, every $1\le p<\infty$, and every
base map $u$, the derivative of the global Jacobian
$J:\dot W^{1,np}\to\mathcal H^p$ is not onto.  Consequently no right
inverse is G\^ateaux differentiable anywhere, and no locally Lipschitz right
inverse exists, even locally over a nonempty open target set.  The argument
includes the exponent $p=2$ omitted by the known general Banach-geometry
obstruction.  Bare surjectivity and merely continuous right inverses remain
open.
}}
\end{center}

\section{Source problem}

The source is Sauli Lindberg, \emph{On the Hardy Space Theory of Compensated
Compactness Quantities}, arXiv:1611.02223, published in 2017 \cite{Lindberg}.
Its Conjecture~7.1, attributed to Iwaniec, states:
\begin{quote}
For every $n\ge2$ and every $p\in[1,\infty)$, the Jacobian operator
\[
 J:\dot W^{1,np}(\R^n,\R^n)\longrightarrow\mathcal H^p(\R^n)
\]
has a continuous right inverse.
\end{quote}
Here $\mathcal H^1$ is the real Hardy space and
$\mathcal H^p=L^p$ with equivalent norms when $p>1$.  The source emphasizes
that the homogeneous scaling removes its inhomogeneous non-surjectivity
argument.

\begin{figure}[htbp]
  \centering
  \includegraphics[width=0.92\textwidth]{figures/open_conjecture_crop.png}
  \caption{Section 7.1 and Conjecture 7.1 in the source PDF, page 25.}
\end{figure}

Later work proves that surjectivity is equivalent to a bounded selection and
to the natural a priori estimate, but the problem remains open
\cite{GKL}.  That work also rules out Fr\'echet-differentiable right inverses
for $p\ne2$ by Banach-space geometry.  The result below is specific to the
Jacobian: it covers every $p$, strengthens Fr\'echet to G\^ateaux
differentiability, and yields a locally Lipschitz no-go theorem.

\section{Proof intuition}

At a map $u$, the Jacobian linearizes to
\[
 DJ(u)v=\cof(Du):Dv.
\]
The coefficient $\cof(Du)$ has a finite global Lebesgue norm.  Translate one
fixed smooth dual test bump into pairwise disjoint regions.  Its dual norm is
unchanged, while the cofactor mass in those regions tends to zero.  Hence the
adjoint of $DJ(u)$ sends unit dual vectors to zero in norm.  The adjoint of a
surjective bounded linear map is bounded below, so $DJ(u)$ cannot be onto.

If a locally Lipschitz section existed, Phelps's infinite-dimensional
Rademacher theorem would make it G\^ateaux differentiable somewhere: the
target of the section is a reflexive homogeneous Sobolev space and therefore
has the Radon--Nikod\'ym property.  Differentiating $J\circ G=\mathrm{Id}$
would force one of the impossible linearizations to be onto.

\section{Statement}

For $1\le p<\infty$, set
\[
 X_p=\dot W^{1,np}(\R^n,\R^n),\qquad
 Y_p=\begin{cases}\mathcal H^1(\R^n),&p=1,\\
                   L^p(\R^n),&p>1.
      \end{cases}
\]
The norm on $X_p$ is $\norm{u}_{X_p}=\norm{Du}_{L^{np}}$, modulo constants.

\begin{theorem}[Nowhere-surjective linearization]\label{thm:linear}
Let $n\ge2$ and $1\le p<\infty$.  For every $u\in X_p$, the Fr\'echet
derivative
\[
 DJ(u):X_p\longrightarrow Y_p
\]
is not surjective.
\end{theorem}

\begin{corollary}[No regular local selection]\label{cor:selection}
Under the hypotheses of Theorem~\ref{thm:linear}:
\begin{enumerate}
\item every set-theoretic right inverse $G:Y_p\to X_p$, if one exists, is
G\^ateaux-differentiable nowhere;
\item there do not exist a nonempty open set $U\subset Y_p$ and a locally
Lipschitz map $G:U\to X_p$ such that $J(G(f))=f$ for every $f\in U$.
\end{enumerate}
\end{corollary}

\section{Proof}

\subsection{The derivative}

The determinant polarization and H\"older's inequality show, for $p>1$,
that $J:X_p\to L^p$ is a continuous $n$-homogeneous polynomial.  At $p=1$,
the multilinear compensated-compactness estimate of Coifman--Lions--Meyer--
Semmes gives the corresponding continuity into $\mathcal H^1$
\cite{CLMS,Lindberg}.  Thus in every case $J$ is Fr\'echet differentiable and
\begin{equation}\label{eq:derivative}
 A_u v:=DJ(u)v=\cof(Du):Dv.
\end{equation}
If
\begin{equation}\label{eq:exponents}
 q=\frac{np}{n-1},\qquad s=(np)'=\frac{np}{np-1},
\end{equation}
then
\[
 C_u:=\cof(Du)\in L^q(\R^n;\R^{n\times n}).
\]

\subsection{Localized dual bumps}

We prove Theorem~\ref{thm:linear}.  First suppose $p>1$.  Choose a nonzero
$\psi\in C_c^\infty(\R^n)$ normalized by
$\norm{\psi}_{L^{p'}}=1$.  Select translations $a_k$ so that
\[
 E_k=a_k+\supp\psi
\]
are pairwise disjoint, and set $\psi_k(x)=\psi(x-a_k)$.  The dual vectors
$\psi_k\in Y_p^*=L^{p'}$ all have norm one.  For $v\in X_p$, equations
\eqref{eq:derivative}--\eqref{eq:exponents} and H\"older's inequality give
\begin{align}
 \big|\langle A_uv,\psi_k\rangle\big|
 &\leq \norm{\psi_k C_u}_{L^s}\norm{Dv}_{L^{np}},\label{eq:dual}\\
 \norm{\psi_k C_u}_{L^s}
 &\leq \norm{\psi_k}_{L^{p'}}
          \norm{1_{E_k}C_u}_{L^q}
  =\norm{1_{E_k}C_u}_{L^q},\label{eq:holder}
\end{align}
because $1/s=1/p'+1/q$.  Since the $E_k$ are disjoint and $C_u\in L^q$,
\begin{equation}\label{eq:vanish}
 \norm{1_{E_k}C_u}_{L^q}\longrightarrow0.
\end{equation}
It follows that
\[
 \norm{A_u^*\psi_k}_{X_p^*}\longrightarrow0
 \quad\text{while}\quad
 \norm{\psi_k}_{Y_p^*}=1.
\]

For $p=1$, choose a nonconstant $\psi\in C_c^\infty(\R^n)$ normalized by
$\norm{\psi}_{\BMO}=1$, where BMO is understood modulo constants, and use
the same disjoint translations.  Translation invariance gives
$\norm{\psi_k}_{(\mathcal H^1)^*}=1$.  Here $q=s=n/(n-1)$, and the compact
bump has finite $L^\infty$ norm, so \eqref{eq:dual} becomes
\[
 \norm{A_u^*\psi_k}_{X_1^*}
 \leq \norm{\psi}_{L^\infty}
       \norm{1_{E_k}C_u}_{L^{n/(n-1)}}\longrightarrow0.
\]
The pairing is the ordinary integral because the test functions are smooth
and compactly supported.

For any bounded surjection $A:X\to Y$, the open mapping theorem implies
that $A^*:Y^*\to X^*$ is bounded below.  The normalized sequences above
contradict that necessary condition.  Therefore $A_u=DJ(u)$ is not onto for
any $u$, proving Theorem~\ref{thm:linear}.

\subsection{Right inverses}

Suppose a right inverse $G$ is G\^ateaux differentiable at $f$.  Since $J$
is Fr\'echet differentiable, the chain rule applied to $J\circ G=\mathrm{Id}$
gives
\[
 DJ(G(f))\,DG(f)=I_{Y_p}.
\]
Thus $DJ(G(f))$ is onto, contrary to Theorem~\ref{thm:linear}.  This proves
the first part of Corollary~\ref{cor:selection}.

For the second part, $Y_p$ is separable.  Moreover, the gradient identifies
$X_p$ isometrically with a closed subspace of $L^{np}$; since $np>1$, $X_p$
is reflexive and hence has the Radon--Nikod\'ym property.  Phelps's theorem
therefore says that a locally Lipschitz map from a nonempty open subset of
$Y_p$ into $X_p$ is G\^ateaux differentiable outside a Gaussian-null set, in
particular at some point \cite{Phelps}.  A locally Lipschitz local section
would contradict the first part.  This completes the proof.

\section{What this settles and what remains open}

The theorem closes the $p=2$ gap in the known \emph{linearization} obstruction
and rules out a broad, natural regularity class of selections.  The proof is
local and requires no assumption that the full Jacobian is onto.

It does not settle the source conjecture.  A merely continuous right inverse
need not be G\^ateaux differentiable anywhere, and the natural scaling of an
$n$-homogeneous equation suggests only $1/n$-H\"older behavior at the origin.
Nor does non-surjectivity of every derivative imply non-surjectivity of a
nonlinear map in infinite dimensions.  Thus bare surjectivity, continuous
selection, the source's sharper energy-minimizer conjecture, and the bounded-
domain variants all remain open.

\paragraph{Verification notes.}
The proof has no computational dependency.  The main checkpoints are the
Hardy-endpoint derivative estimate, the exponent identity
$1/(np)'=1/p'+(n-1)/(np)$, the use of pairwise disjoint translated supports,
and the RNP hypothesis in Phelps's theorem.

\paragraph{Novelty check.}
On 11 August 2026 we searched the run's four lightweight indexes and the exact
source/later-paper titles, together with combinations of ``linearized
Jacobian not surjective'', ``Jacobian submersion nowhere $p=2$'', ``locally
Lipschitz right inverse Jacobian'', and ``Phelps Jacobian right inverse''.
The closest result found is the general Banach-geometry obstruction in
\cite{GKL}: no submersion for $p<2$, no regular submersion for $p>2$, and no
Fr\'echet-differentiable right inverse for $p\ne2$.  No source located the
cofactor-localization proof, the all-$p$ statement, or the locally Lipschitz
corollary.  Novelty confidence is moderate, not definitive.

\paragraph{Human-review recommendation.}
Review as a substantial partial result.  In addition to the four checkpoints
above, compare carefully with the published version of \cite{GKL} to confirm
that the Jacobian-specific $p=2$ observation and Phelps upgrade have not been
recorded elsewhere.

\begin{thebibliography}{9}
\bibitem{Lindberg}
S.~Lindberg,
\emph{On the Hardy Space Theory of Compensated Compactness Quantities},
Arch. Ration. Mech. Anal. \textbf{224} (2017), 709--742;
arXiv:1611.02223.

\bibitem{CLMS}
R.~Coifman, P.-L.~Lions, Y.~Meyer, and S.~Semmes,
\emph{Compensated Compactness and Hardy Spaces},
J. Math. Pures Appl. \textbf{72} (1993), 247--286.

\bibitem{GKL}
A.~Guerra, L.~Koch, and S.~Lindberg,
\emph{Nonlinear Open Mapping Principles, with Applications to the Jacobian
Equation and Other Scale-Invariant PDEs},
arXiv:2010.10497 (2020).

\bibitem{Phelps}
R.~R. Phelps,
\emph{Gaussian Null Sets and Differentiability of Lipschitz Maps on Banach
Spaces},
Pacific J. Math. \textbf{77} (1978), 523--531.
\end{thebibliography}

\end{document}
